\documentclass[11pt]{article}

\usepackage[utf8]{inputenc}
\usepackage[T1]{fontenc}
\usepackage[spanish,es-nodecimaldot]{babel}
\usepackage[margin=1in]{geometry}
\usepackage{graphicx}
\usepackage{booktabs}
\usepackage{longtable}
\usepackage{amsmath}
\usepackage{amssymb}
\usepackage{hyperref}
\usepackage{float}
\usepackage{array}

\hypersetup{
    colorlinks=true,
    linkcolor=black,
    citecolor=black,
    urlcolor=black
}

\newcommand{\figroot}{../../outputs/observational_bias_audit/figures}
\newcommand{\tabroot}{../../outputs/observational_bias_audit/tables}

% The following wrappers allow the document to compile even on machines where
% the generated figures or tables have not been copied yet. When the project
% outputs exist, the real files are inserted normally.
\newcommand{\safeincludegraphics}[2][]{%
  \IfFileExists{#2}{%
    \includegraphics[#1]{#2}%
  }{%
    \fbox{\begin{minipage}{0.86\textwidth}\centering
    Archivo de figura no encontrado:\\[0.5em]
    \texttt{\detokenize{#2}}
    \end{minipage}}%
  }%
}

\newcommand{\safetableinput}[1]{%
  \IfFileExists{#1}{%
    \input{#1}%
  }{%
    \begin{center}
    \fbox{\begin{minipage}{0.86\textwidth}\centering
    Archivo de tabla no encontrado:\\[0.5em]
    \texttt{\detokenize{#1}}
    \end{minipage}}
    \end{center}
  }%
}

\title{Auditoría observacional del Mapper orbital en el catálogo de exoplanetas}
\author{}
\date{}

\begin{document}
\maketitle

\section{Introducción}

Este subproyecto evalúa si la topología observada en Mapper, en particular en el espacio orbital, está asociada con el método de descubrimiento de los planetas. La pregunta científica no es todavía si Mapper descubre clases planetarias definitivas, sino si la estructura del grafo contiene una señal observacional suficientemente fuerte como para condicionar cualquier interpretación física posterior.

En términos prácticos, este análisis funciona como una auditoría post-hoc: el grafo Mapper y sus membresías nodo-planeta se mantienen fijos, y se agregan etiquetas externas de descubrimiento para medir si los nodos, componentes y regiones periféricas están enriquecidos por métodos como tránsito, velocidad radial, imagen directa o microlensing.

\section{Hipótesis}

La hipótesis evaluada es:
\[
H_A: \quad \text{la estructura topológica de Mapper está explicada de forma importante por sesgos de observación y método de descubrimiento.}
\]

La hipótesis rival es:
\[
H_F: \quad \text{la estructura topológica conserva una señal física u orbital no reducible al método de descubrimiento.}
\]

Este documento no prueba causalidad. Una asociación estadísticamente fuerte entre nodo Mapper y método de descubrimiento apoya que el sesgo observacional es una componente relevante de la topología, pero no demuestra por sí sola que sea la causa principal. Para afirmar causalidad o dominancia causal se requiere una segunda vía contrafactual: balanceo por método, residualización o reconstrucción de Mapper bajo controles observacionales.

\section{Datos y configuración analizada}

El análisis principal usa la configuración \texttt{orbital\_pca2\_cubes10\_overlap0p35}. Esta decisión es metodológicamente razonable porque el Mapper orbital ya había sido identificado como el caso con mejor equilibrio entre complejidad topológica y baja dependencia de imputación. Como comparación se resumen también, cuando existen artifacts disponibles, las configuraciones \texttt{phys\_min\_pca2\_cubes10\_overlap0p35}, \texttt{joint\_no\_density\_pca2\_cubes10\_overlap0p35}, \texttt{joint\_pca2\_cubes10\_overlap0p35} y \texttt{thermal\_pca2\_cubes10\_overlap0p35}.

La unidad estadística de este subproyecto no es solamente el planeta individual, sino la incidencia planeta--nodo inducida por Mapper. Esto importa porque un planeta puede aparecer en más de un nodo debido al solapamiento de cubiertas propio del algoritmo Mapper.

\section{Metodología}

Sea $G=(V,E)$ el grafo Mapper, donde $V$ es el conjunto de nodos y $E$ el conjunto de aristas. Cada nodo $v\in V$ contiene un subconjunto de planetas $S_v$. Para cada planeta $i\in S_v$, sea $d_i$ su método de descubrimiento. Para cada nodo $v$ y método $m$, se calcula:
\[
p_v(m) = \frac{\mathrm{count}_v(m)}{n_v},
\]
con $n_v=|S_v|$.

La pureza nodal se define como:
\[
\mathrm{purity}(v) = \max_m p_v(m).
\]

La entropía de métodos dentro del nodo se define como:
\[
H(v) = -\sum_m p_v(m)\log p_v(m),
\]
con la convención $0\log(0)=0$. La entropía normalizada es:
\[
H_{\mathrm{norm}}(v) = \frac{H(v)}{\log(K)},
\]
where $K$ es el número de métodos de descubrimiento presentes en el universo analizado.

Interpretación:
\[
\mathrm{purity}(v)\approx 1 \quad \Rightarrow \quad \text{nodo dominado por un método de descubrimiento},
\]
\[
H_{\mathrm{norm}}(v)\approx 0 \quad \Rightarrow \quad \text{baja mezcla metodológica},
\]
\[
H_{\mathrm{norm}}(v)\text{ alto} \quad \Rightarrow \quad \text{nodo observacionalmente heterogéneo}.
\]

La prueba nula mantiene fijo el grafo y la membresía nodo-planeta, permuta \texttt{discoverymethod} entre planetas y recalcula las métricas globales. El z-score reportado sigue:
\[
z = \frac{\mathrm{observed}-\mathrm{null\_mean}}{\mathrm{null\_std}}.
\]

El valor-$p$ empírico se interpreta como evidencia contra la hipótesis nula de que la asociación nodo--método surge por asignación aleatoria de métodos al grafo fijo.

\section{Métricas nodo--método}

Se generan tablas de pureza, entropía, método dominante, tamaño nodal, fracción de imputación media, fracción de variables físicamente derivadas, grado, componente e indicador de periferia. Estas métricas permiten separar tres escenarios:

\begin{enumerate}
    \item \textbf{Sesgo observacional localizado}: pocos nodos o componentes están enriquecidos por método.
    \item \textbf{Sesgo observacional global}: gran parte del grafo tiene baja entropía y alta pureza por método.
    \item \textbf{Mezcla física-observacional}: el núcleo del grafo mezcla métodos, pero las ramas o extremos muestran enriquecimientos particulares.
\end{enumerate}

\section{Prueba nula por permutación}

Las métricas globales contrastadas contra la distribución nula son pureza media ponderada, entropía media ponderada, información mutua normalizada entre incidencias nodo--método y assortativity del método dominante entre nodos conectados cuando existen aristas.

La métrica más importante para este documento es la información mutua normalizada entre nodo y método, denotada aquí como NMI. Si la NMI observada cae muy lejos de la distribución nula, entonces la asignación de métodos de descubrimiento no es intercambiable respecto a la topología Mapper: el grafo conserva información sobre cómo fueron descubiertos los planetas.

\section{Resultados e interpretación de figuras}

\subsection{Método dominante en el grafo orbital}

\begin{figure}[H]
\centering
\safeincludegraphics[width=0.9\textwidth]{\figroot/orbital_graph_by_dominant_discovery_method.pdf}
\caption{Grafo Mapper orbital coloreado por método de descubrimiento dominante.}
\label{fig:dominant-method}
\end{figure}

\paragraph{Interpretación.}
La Figura~\ref{fig:dominant-method} muestra que el grafo orbital no está coloreado de forma homogénea. La mayor parte de los nodos grandes y del componente principal están dominados por tránsito o velocidad radial, mientras que métodos menos frecuentes, como imagen directa o microlensing, aparecen en regiones más localizadas. Esto sugiere que la estructura topológica no debe leerse inmediatamente como una taxonomía física de planetas: una parte de la separación nodal puede estar inducida por las ventanas de sensibilidad de cada técnica de detección.

\paragraph{Implicación.}
El resultado apoya una versión moderada de la Hipótesis A: el método de descubrimiento sí está codificado en el grafo. Sin embargo, esta figura por sí sola no demuestra que el método sea la causa principal de toda la topología; solamente muestra una asociación espacial clara entre regiones Mapper y etiquetas observacionales.

\subsection{Pureza nodal por método dominante}

\begin{figure}[H]
\centering
\safeincludegraphics[width=0.9\textwidth]{\figroot/orbital_graph_by_method_purity.pdf}
\caption{Pureza nodal por método de descubrimiento. Valores cercanos a 1 indican nodos dominados por un solo método.}
\label{fig:purity}
\end{figure}

\paragraph{Interpretación.}
La Figura~\ref{fig:purity} indica que muchos nodos tienen pureza alta, con valores cercanos a 1. Esto significa que varios nodos contienen planetas descubiertos casi exclusivamente por un mismo método. El patrón no se limita visualmente a nodos aislados: también aparecen nodos relativamente grandes con pureza elevada. Por lo tanto, la señal observacional no parece ser únicamente un artefacto de nodos diminutos.

\paragraph{Lectura científica.}
Una pureza alta es esperable si ciertos métodos exploran regiones distintas del espacio orbital. Por ejemplo, tránsito favorece detecciones de periodos cortos y geometrías alineadas; velocidad radial favorece planetas con señales dinámicas detectables; imagen directa favorece objetos más masivos y separados. Si estas preferencias entran en las variables orbitales usadas por Mapper, el grafo puede separar parcialmente métodos antes de separar clases físicas intrínsecas.

\subsection{Entropía normalizada nodo--método}

\begin{figure}[H]
\centering
\safeincludegraphics[width=0.9\textwidth]{\figroot/orbital_graph_by_method_entropy.pdf}
\caption{Entropía normalizada nodo--método en el Mapper orbital. Valores bajos indican nodos dominados por pocos métodos; valores altos indican mezcla observacional.}
\label{fig:entropy}
\end{figure}

\paragraph{Interpretación.}
La Figura~\ref{fig:entropy} complementa la lectura de pureza. La presencia de numerosos nodos con entropía baja indica que muchos nodos no mezclan uniformemente los métodos de descubrimiento. Al mismo tiempo, algunos nodos del componente central tienen entropía mayor, lo que sugiere que el núcleo del grafo puede funcionar como una región de mezcla entre poblaciones observacionales.

\paragraph{Conclusión local.}
La topología orbital parece contener una combinación de regiones homogéneas por método y zonas de transición. Esta estructura es consistente con una lectura mixta: hay una señal observacional fuerte, pero no necesariamente una partición absoluta del catálogo por método.

\subsection{Composición método--nodo en los nodos más grandes}

\begin{figure}[H]
\centering
\safeincludegraphics[width=0.9\textwidth]{\figroot/orbital_node_method_fraction_heatmap.pdf}
\caption{Heatmap nodo x método para el espacio orbital. Se muestran los nodos más grandes o más relevantes según el criterio de generación de la figura.}
\label{fig:heatmap}
\end{figure}

\paragraph{Interpretación.}
La Figura~\ref{fig:heatmap} muestra que los nodos grandes están dominados principalmente por tránsito y velocidad radial. Los métodos minoritarios aparecen como contribuciones más escasas y localizadas. Esta distribución es importante porque evita una interpretación ingenua: el sesgo observacional no solo aparece en métodos raros o nodos pequeños, sino en la estructura dominante del catálogo, donde tránsito y velocidad radial explican gran parte de la masa estadística.

\paragraph{Implicación.}
La lectura física de nodos grandes debe condicionarse al método dominante. Un nodo dominado por tránsito puede estar capturando un régimen real de planetas compactos de periodo corto, pero también puede reflejar la región del espacio de parámetros donde el método de tránsito es más eficiente.

\subsection{Prueba nula de NMI}

\begin{figure}[H]
\centering
\safeincludegraphics[width=0.75\textwidth]{\figroot/orbital_permutation_null_nmi.pdf}
\caption{Distribución nula de NMI y valor observado. La línea vertical representa la NMI observada.}
\label{fig:null-nmi}
\end{figure}

\paragraph{Interpretación.}
La Figura~\ref{fig:null-nmi} es la evidencia estadística más fuerte de este subproyecto. La distribución nula queda concentrada cerca de valores pequeños de NMI, mientras que el valor observado aparece muy desplazado hacia la derecha. En la tabla global, el Mapper orbital tiene NMI observada de aproximadamente $0.165$, z-score global de aproximadamente $138.9$ y valor-$p$ empírico de $0.001$.

\paragraph{Conclusión estadística.}
Bajo el diseño de permutación usado, la asociación entre nodo Mapper y método de descubrimiento es mucho mayor que la esperada por azar. Por tanto, se rechaza la explicación de que la distribución de métodos sobre el grafo orbital sea aleatoria respecto a la topología fija.

\subsection{Pureza frente a imputación}

\begin{figure}[H]
\centering
\safeincludegraphics[width=0.75\textwidth]{\figroot/orbital_purity_vs_imputation.pdf}
\caption{Pureza nodal frente a fracción media de imputación.}
\label{fig:purity-imputation}
\end{figure}

\paragraph{Interpretación.}
La Figura~\ref{fig:purity-imputation} muestra que muchos nodos de pureza alta tienen fracción media de imputación muy baja. Esto es relevante porque reduce la posibilidad de que la asociación nodo--método sea solamente un artefacto de valores completados. En particular, el resumen global reporta para el Mapper orbital una imputación media cercana a $0.011$, muy inferior a la del espacio térmico.

\paragraph{Implicación.}
En el caso orbital, la evidencia de sesgo observacional no parece depender principalmente de imputación masiva. Esto hace que la señal sea más interesante: el grafo puede estar capturando un sesgo real del proceso de observación, no solo ruido introducido por el pipeline de completación de datos.

\subsection{Pureza frente a tamaño de nodo}

\begin{figure}[H]
\centering
\safeincludegraphics[width=0.75\textwidth]{\figroot/orbital_purity_vs_node_size.pdf}
\caption{Pureza nodal frente a tamaño de nodo.}
\label{fig:purity-size}
\end{figure}

\paragraph{Interpretación.}
La Figura~\ref{fig:purity-size} permite controlar visualmente un problema estadístico importante: los nodos pequeños tienden a ser puros por azar o por baja diversidad muestral. Ese efecto sí aparece, pero no explica todo el patrón, porque también existen nodos grandes con pureza elevada. La presencia de nodos masivos dominados por un método indica que la señal observacional opera también en regiones poblacionalmente importantes del grafo.

\paragraph{Matiz.}
La figura también muestra algunos nodos grandes con pureza más baja, lo que sugiere que no todo el núcleo está dominado por un único método. Por tanto, la conclusión correcta no es ``todo el Mapper es sesgo'', sino que la topología contiene una señal observacional fuerte y heterogénea.

\subsection{Composición por componente conexa}

\begin{figure}[H]
\centering
\safeincludegraphics[width=0.9\textwidth]{\figroot/orbital_component_method_composition.pdf}
\caption{Composición por método de descubrimiento a nivel de componente.}
\label{fig:component-composition}
\end{figure}

\paragraph{Interpretación.}
La Figura~\ref{fig:component-composition} muestra que las componentes conexas no tienen la misma mezcla de métodos. Algunas están casi completamente dominadas por tránsito, mientras que otras concentran mayor proporción de velocidad radial o métodos minoritarios. Esta evidencia es más estructural que la pureza de nodos individuales, porque indica que la segmentación por método persiste a escala de componente.

\paragraph{Implicación.}
Si una componente completa está dominada por un método, entonces las ramas o islas de Mapper pueden estar reflejando la geometría observacional del catálogo. Esto es especialmente importante para componentes pequeñas, donde la interpretación física debe hacerse con cautela.

\subsection{Comparación entre configuraciones}

\begin{figure}[H]
\centering
\safeincludegraphics[width=0.95\textwidth]{\figroot/config_comparison_bias_metrics.pdf}
\caption{Comparación resumida de métricas de sesgo observacional entre configuraciones disponibles.}
\label{fig:config-comparison}
\end{figure}

\paragraph{Interpretación.}
La Figura~\ref{fig:config-comparison} muestra que la asociación entre topología y método de descubrimiento no es exclusiva del Mapper orbital. Las configuraciones conjuntas también presentan NMI y z-scores elevados. Sin embargo, la interpretación no debe basarse solo en fuerza de asociación: también importa la imputación. El espacio térmico presenta una imputación media mucho mayor, por lo que cualquier sesgo observado allí es más difícil de separar de efectos del pipeline de completación.

\paragraph{Lectura comparativa.}
El caso orbital sigue siendo el más útil para discutir sesgo observacional porque combina tres propiedades: alta señal nodo--método, baja imputación media y complejidad topológica previamente identificada como relevante. Las configuraciones conjuntas pueden amplificar la asociación con método, pero también introducen mayor mezcla entre variables físicas, orbitales y derivadas.

\section{Resultados cuantitativos principales}

La Tabla~\ref{tab:global-summary-ref} resume los valores globales. En el Mapper orbital se observa una pureza media ponderada aproximada de $0.8422$, una entropía media ponderada aproximada de $0.4370$, una entropía normalizada de $0.1898$, una NMI nodo--método de $0.1650$, assortativity del método dominante de $0.8018$ y un valor-$p$ empírico de $0.001$ para la NMI global. Estos valores apuntan a una asociación fuerte entre estructura Mapper y método de descubrimiento.

Un punto importante es la comparación central--periferia. Para el Mapper orbital, la tabla reporta $79$ nodos centrales y $45$ periféricos. Los nodos centrales tienen pureza media aproximada de $0.8881$, entropía media de $0.2625$ e imputación media de $0.0054$; los nodos periféricos tienen pureza media aproximada de $0.8159$, entropía media de $0.3986$ e imputación media de $0.0215$. Por tanto, en esta corrida la asociación con método no debe describirse como exclusivamente periférica. Hay evidencia de sesgo tanto en el núcleo como en regiones periféricas, aunque los métodos minoritarios tienden a localizarse en zonas más específicas.

\section{Discusión}

Los resultados apoyan que el método de descubrimiento está estadísticamente acoplado con la topología Mapper. Esto es consistente con la astronomía observacional: cada técnica de detección induce una función de selección distinta sobre periodo orbital, separación, masa, radio, inclinación y brillo estelar. Por eso, aun si las variables usadas por Mapper son físicas u orbitales, el grafo puede heredar la geometría del proceso de descubrimiento.

La conclusión más defendible no es que la interpretación física sea falsa. La conclusión correcta es más fina:
\[
\text{topología observada} = \text{señal física/orbital} + \text{sesgo observacional} + \text{efecto de imputación} + \text{sensibilidad Mapper}.
\]

En el espacio orbital, la baja imputación media hace que la señal de sesgo observacional sea especialmente relevante: no parece ser solamente un artefacto de completación de datos. Sin embargo, el análisis sigue siendo post-hoc. Para afirmar que el sesgo es la causa principal de la topología, falta una segunda vía: construir Mappers balanceados por método, residualizar variables contra método de descubrimiento o comparar grafos reconstruidos bajo controles observacionales.

\section{Conclusiones}

\begin{enumerate}
    \item \textbf{La Hipótesis A recibe apoyo correlacional fuerte.} La NMI observada entre nodo Mapper y método de descubrimiento es mucho mayor que la esperada bajo permutación aleatoria, con valor-$p$ empírico de $0.001$ en la configuración orbital.
    \item \textbf{El sesgo observacional no es marginal.} Muchos nodos tienen alta pureza por método y baja entropía, lo que indica que la topología contiene información sobre el proceso de descubrimiento.
    \item \textbf{El resultado orbital es especialmente importante porque tiene baja imputación.} La asociación nodo--método aparece en un espacio donde la fracción media de imputación es baja, lo que hace menos probable que el patrón sea solo un artefacto de datos completados.
    \item \textbf{No se debe concluir todavía que no existen clases físicas.} El análisis muestra que el método de descubrimiento explica una parte significativa de la estructura, pero no cuantifica cuánta estructura queda después de controlar ese método.
    \item \textbf{La siguiente prueba necesaria es contrafactual.} El siguiente subproyecto debe construir versiones balanceadas o residualizadas del Mapper para medir si la topología sobrevive cuando se controla explícitamente el método de descubrimiento.
\end{enumerate}

\section{Limitaciones}

Este análisis depende de la calidad de la metadata observacional, del esquema de imputación ya fijado en el pipeline original y de membresías nodales con solapamiento propio de Mapper. Además, el método de descubrimiento no es una variable causal aislada: está correlacionado con época de descubrimiento, instrumento, misión, precisión fotométrica o espectroscópica, tipo de estrella, brillo aparente y criterios de publicación.

Otra limitación es que la tabla de nodos enriquecidos puede estar dominada por configuraciones no orbitales si se reportan todas las configuraciones juntas. Para una discusión estrictamente orbital conviene generar también una tabla filtrada, por ejemplo \texttt{top\_enriched\_nodes\_orbital.tex}, con los nodos enriquecidos únicamente de \texttt{orbital\_pca2\_cubes10\_overlap0p35}.

\section{Conclusión general}

El subproyecto cuantifica una asociación fuerte entre topología Mapper y método de descubrimiento. Esto no invalida el análisis físico, pero obliga a reinterpretarlo: cualquier región candidata del Mapper orbital debe reportarse junto con su composición observacional, su pureza por método, su entropía, su imputación media y su estabilidad bajo controles posteriores.

La conclusión sintética es:
\[
\boxed{
\text{El Mapper orbital contiene señal topológica real, pero esa señal está observacionalmente estratificada.}
}
\]

\appendix
\section{Apéndice de tablas}

\subsection{Resumen global de métricas de sesgo observacional}
\label{tab:global-summary-ref}
\safetableinput{\tabroot/summary_global_bias_metrics.tex}

\subsection{Nodos con mayor enriquecimiento por método de descubrimiento}
\safetableinput{\tabroot/top_enriched_nodes.tex}

\subsection{Comparación entre nodos centrales y periféricos}
\safetableinput{\tabroot/central_vs_peripheral_bias.tex}

\end{document}
